%%%%%%%%%%%%%%%%%%%%%%preface.tex%%%%%%%%%%%%%%%%%%%%%%%%%%%%%%%%%%%%%%%%%
% sample preface
%
% Use this file as a template for your own input.
%
%%%%%%%%%%%%%%%%%%%%%%%% Springer Nature %%%%%%%%%%%%%%%%%%%%%%%%

\preface

The quest for a unified theory of artificial intelligence (AI) has been a long-standing goal in the field of AI research. The idea of combining symbolic reasoning with neural networks, known as neuro-symbolic AI, has gained significant attention in recent years as a promising approach to achieving this goal.
This book explores the tensor network approach to neuro-symbolic AI, which is a powerful mathematical framework for representing and manipulating complex data structures.



Use the template \emph{preface.tex} together with the document class \verb|SNmono| (monograph-type books) or \verb|SNmult| (edited books) to style your preface.

A preface\index{preface} is a book's preliminary statement, usually written by the \textit{author or editor} of a work, which states its origin, scope, purpose, plan, and intended audience, and which sometimes includes afterthoughts and acknowledgments of assistance. 

When written by a person other than the author, it is called a foreword.
The preface or foreword is distinct from the introduction, which deals with the subject of the work.

Customarily \textit{acknowledgments} are included as last part of the preface.
 

\vspace{\baselineskip}
\begin{flushright}\noindent
Place(s),\hfill {\it Firstname  Surname}\\
month year%\hfill {\it Firstname  Surname}\\
\end{flushright}


